% ==============================================================================
% CAPITOLO 11: VERIFICA DELLE IPOTESI STATISTICHE (TEST DI SIGNIFICATIVITÀ)
% ==============================================================================
\chapter{Verifica delle Ipotesi Statistiche (Test di Significatività)}
\label{cap:test_statistici}

\begin{abstract}
Il presente capitolo espone la teoria formale delle decisioni statistiche di Jerzy Neyman ed Egon Pearson. Si formalizza la dialettica tra ipotesi nulla ($H_0$) e ipotesi alternativa ($H_1$), chiarendo l'asimmetria ontologica tra le due posizioni mediante la celebre metafora giuridica del processo penale. Si analizza la matrice di contingenza degli errori decisionali (Errore di I tipo $\alpha$ ed Errore di II tipo $\beta$), definendo la funzione di potenza del test. Viene approfondito il concetto unificante di $p$-value e dimostrato il teorema di dualità strutturale che lega i test bilaterali agli intervalli di confidenza. Infine, si presenta la tassonomia operativa dei test parametrici $Z$ e $t$ per medie, proporzioni e confronti tra campioni indipendenti.
\end{abstract}

% ==============================================================================
% SEZIONE 11.1: STRUTTURA LOGICA E METAFORA GIURIDICA DI NEYMAN-PEARSON
% ==============================================================================
\section{La Struttura Logica del Test d'Ipotesi e la Metafora Giuridica}
\label{sec:struttura_logica_test}

Nel metodo scientifico e nell'ingegneria gestionale, il decisore deve costantemente vagliare affermazioni quantitative su parametri ignoti: un nuovo trattamento riduce i consumi? Il tasso di difettosità supera la soglia di tolleranza contrattuale? 

La teoria della verifica delle ipotesi di Neyman-Pearson non cerca di ``dimostrare la verità assoluta'', bensì stabilisce una procedura oggettiva per decidere se l'evidenza campionaria è sufficientemente forte da respingere una teoria di riferimento.

\begin{definizione}{Ipotesi Nulla ($H_0$) e Ipotesi Alternativa ($H_1$)}{def_ipotesi_statistica}
\begin{itemize}
    \item \textbf{Ipotesi Nulla ($H_0$):} Rappresenta lo \emph{status quo}, l'ipotesi di neutralità, assenza di effetto o conformità nominale (ad esempio $H_0: \mu = \mu_0$). Viene assunta provvisoriamente vera fino a prova contraria (principio di conservazione).
    \item \textbf{Ipotesi Alternativa ($H_1$ o $H_a$):} Rappresenta la congettura del ricercatore o l'effetto innovativo che si intende sostenere mediante evidenza empirica schiacciante:
    \begin{itemize}
        \item \textbf{Bilaterale (a due code):} $H_1: \mu \neq \mu_0$;
        \item \textbf{Unilaterale destra (a una coda):} $H_1: \mu > \mu_0$;
        \item \textbf{Unilaterale sinistra (a una coda):} $H_1: \mu < \mu_0$.
    \end{itemize}
\end{itemize}
\end{definizione}

\paragraph{La Metafora Giuridica del Processo Penale:}
Per comprendere la profonda asimmetria tra $H_0$ e $H_1$:
\begin{itemize}
    \item $H_0$ equivale alla **presunzione di innocenza dell'imputato** (innocente fino a prova contraria oltre ogni ragionevole dubbio).
    \item $H_1$ equivale alla **colpevolezza**.
    \item Il verdetto del tribunale non è mai ``dimostrazione di innocenza'', bensì:
    \begin{itemize}
        \item \textbf{Rifiuto di $H_0$ (Condanna):} Le prove empiriche a carico sono talmente schiaccianti che l'ipotesi di innocenza è incompatibile con i fatti.
        \item \textbf{Non rifiuto di $H_0$ (Assoluzione per insufficienza di prove):} Non significa aver provato che l'imputato è innocente, ma che i dati raccolti non sono sufficienti a condannarlo.
    \end{itemize}
\end{itemize}

% ==============================================================================
% SEZIONE 11.2: MATRICE DECISIONALE, ERRORI E FUNZIONE DI POTENZA
% ==============================================================================
\section{Matrice degli Errori Decisionali e Potenza del Test}
\label{sec:matrice_errori_potenza}

Poiché la decisione è assunta sulla base di un campione casuale limitato, il test è esposto a due tipologie di errore probabilistico:

\begin{table}[htbp]
    \centering
    \caption{Matrice di contingenza e probabilità degli esiti decisionali nel test d'ipotesi.}
    \label{tab:matrice_errori}
    \begin{tabular}{lcc}
        \toprule
        \textbf{Decisione Assunta} & $\bm{H_0}$ \textbf{Vera nella Realtà} & $\bm{H_0}$ \textbf{Falsa nella Realtà ($H_1$ vera)} \\
        \midrule
        \textbf{Non Rifiutare $H_0$} & Decisione Corretta ($1 - \alpha$) & \textbf{Errore di II Tipo} ($\beta$) (\emph{Falso Negativo}) \\[8pt]
        \textbf{Rifiutare $H_0$} & \textbf{Errore di I Tipo} ($\alpha$) (\emph{Falso Positivo}) & Decisione Corretta / \textbf{Potenza} ($1 - \beta$) \\
        \bottomrule
    \end{tabular}
\end{table}

\begin{definizione}{Livello di Significatività ed Errore di I Tipo}{def_livello_significativita}
L'\textbf{Errore di I Tipo} consiste nel condannare un innocente (rifiutare $H_0$ quando $H_0$ è vera). La massima probabilità tollerata per questo errore è fissata a priori dal ricercatore come \textbf{livello di significatività} $\alpha$:
\begin{equation}
    \alpha \coloneqq \mathbb{P}(\text{Rifiutare } H_0 \mid H_0 \text{ è vera}) \qquad (\text{tipicamente } 0.05 \text{ o } 0.01)
\end{equation}
\end{definizione}

\begin{definizione}{Errore di II Tipo e Potenza del Test ($1-\beta$)}{def_potenza_test}
L'\textbf{Errore di II Tipo} consiste nell'assolvere un colpevole (non rifiutare $H_0$ quando nella realtà $H_0$ è falsa). La sua probabilità è $\beta = \mathbb{P}(\text{Non Rifiutare } H_0 \mid H_1 \text{ vera})$.
La \textbf{Potenza del Test} (Power) è la probabilità di rifiutare correttamente $H_0$ quando l'effetto reale esiste:
\begin{equation}
    \text{Potenza}(\theta) \coloneqq 1 - \beta(\theta) = \mathbb{P}(\text{Rifiutare } H_0 \mid \theta \in H_1)
\end{equation}
\end{definizione}

\begin{figure}[htbp]
    \centering
    \begin{tikzpicture}[scale=0.85, >=Stealth]
    \draw[->, thick, navyblue] (-4.5,0) -- (6.5,0) node[right, font=\footnotesize] {$z$};

    % Curva sotto H0 (centrata in 0)
    \draw[domain=-4.0:4.0, smooth, variable=\x, very thick, coolblue]
        plot ({\x}, {3.2*exp(-\x*\x/2)});
    \node[above, font=\small\bfseries, text=coolblue] at (0,3.2) {Sotto $H_0$ ($\mu = \mu_0$)};

    % Curva sotto H1 (centrata in 2.8)
    \draw[domain=-1.0:6.5, smooth, variable=\x, very thick, crimsonred]
        plot ({\x}, {3.2*exp(-(\x-2.8)*(\x-2.8)/2)});
    \node[above, font=\small\bfseries, text=crimsonred] at (2.8,3.2) {Sotto $H_1$ ($\mu = \mu_1$)};

    % Soglia critica z_crit = 1.64
    \draw[ultra thick, dashed, black] (1.64,0) -- (1.64,3.5) node[above, font=\scriptsize\bfseries] {Soglia Critica $z_\alpha$};

    % Errore Tipo I (alpha) sotto H0 a destra di 1.64
    \fill[coolblue!50, domain=1.64:4.0, variable=\x]
        (1.64,0) -- plot ({\x}, {3.2*exp(-\x*\x/2)}) -- (4.0,0) -- cycle;
    \node[font=\tiny\bfseries, text=coolblue!90!black] at (2.3,0.3) {$\alpha$ (Tipo I)};

    % Errore Tipo II (beta) sotto H1 a sinistra di 1.64
    \fill[crimsonred!35, domain=-0.5:1.64, variable=\x]
        (-0.5,0) -- plot ({\x}, {3.2*exp(-(\x-2.8)*(\x-2.8)/2)}) -- (1.64,0) -- cycle;
    \node[font=\tiny\bfseries, text=crimsonred!90!black] at (1.0,0.5) {$\beta$ (Tipo II)};

    % Potenza (1 - beta)
    \node[font=\scriptsize\bfseries, text=forestgreen!80!black] at (3.8,1.2) {Potenza $= 1 - \beta$};
\end{tikzpicture}

    \caption{Sovrapposizione delle distribuzioni campionarie della statistica test sotto $H_0$ (distribuzione blu a sinistra) e sotto $H_1$ per un valore alternativo $\theta_1 > \theta_0$ (distribuzione rossa a destra). Fissata la soglia critica $c_{\alpha}$, l'area rossa a destra di $c_{\alpha}$ sotto la curva $H_0$ è $\alpha$; l'area verde tratteggiata sotto la curva $H_1$ a sinistra di $c_{\alpha}$ è $\beta$; l'area a destra sotto la curva $H_1$ rappresenta la Potenza $1-\beta$.}
    \label{fig:errori_tipo1_tipo2}
\end{figure}

% ==============================================================================
% SEZIONE 11.3: REGIONI DI ACCETTAZIONE E RIFIUTO
% ==============================================================================
\section{Regioni di Accettazione, Regioni Critiche di Rifiuto e Algoritmo Operativo}
\label{sec:regioni_rifiuto}

\begin{figure}[htbp]
    \centering
    \begin{tikzpicture}[scale=0.85, >=Stealth]
    \draw[->, thick, navyblue] (-4.5,0) -- (4.5,0) node[right, font=\footnotesize] {$z$};
    \draw[->, thick, navyblue] (0,-0.2) -- (0,3.8) node[above, font=\footnotesize] {$f(z \mid H_0)$};

    % Regione Accettazione
    \fill[forestgreen!20, domain=-1.96:1.96, variable=\x]
        (-1.96,0) -- plot ({\x}, {3.5*exp(-\x*\x/2)}) -- (1.96,0) -- cycle;

    % Regioni Rifiuto
    \fill[crimsonred!40, domain=-4.0:-1.96, variable=\x]
        (-4.0,0) -- plot ({\x}, {3.5*exp(-\x*\x/2)}) -- (-1.96,0) -- cycle;
    \fill[crimsonred!40, domain=1.96:4.0, variable=\x]
        (1.96,0) -- plot ({\x}, {3.5*exp(-\x*\x/2)}) -- (4.0,0) -- cycle;

    \draw[domain=-4.0:4.0, smooth, variable=\x, very thick, navyblue]
        plot ({\x}, {3.5*exp(-\x*\x/2)});

    \draw[thick, dashed, crimsonred] (-1.96,0) -- (-1.96, {3.5*exp(-1.96*1.96/2)});
    \draw[thick, dashed, crimsonred] (1.96,0) -- (1.96, {3.5*exp(-1.96*1.96/2)});

    \node[below, font=\scriptsize\bfseries, text=crimsonred] at (-1.96,0) {$-z_{\alpha/2}$};
    \node[below, font=\scriptsize\bfseries, text=crimsonred] at (1.96,0) {$+z_{\alpha/2}$};

    \node[font=\bfseries\small, text=forestgreen!80!black] at (0,1.2) {Regione di Accettazione ($1-\alpha$)};
    \node[font=\tiny\bfseries, text=crimsonred, align=center] at (-2.8,0.8) {Rifiuto $H_0$\\($\alpha/2$)};
    \node[font=\tiny\bfseries, text=crimsonred, align=center] at (2.8,0.8) {Rifiuto $H_0$\\($\alpha/2$)};
\end{tikzpicture}

    \caption{Suddivisione dello spazio della statistica test gaussiana sotto $H_0$ per un test bilaterale al livello $\alpha = 5\%$. La regione centrale di Non Rifiuto ha probabilità $1-\alpha = 95\%$, delimitata dai valori critici $\pm z_{0.025} = \pm 1.96$. Le due code estreme rosse costituiscono la Regione Critica di Rifiuto $\mathcal{R}$.}
    \label{fig:test_regioni_rifiuto}
\end{figure}

\begin{metodo}[Procedura Operativa in 4 Fasi per l'Esecuzione di un Test d'Ipotesi]
\begin{enumerate}
    \item \textbf{Formulazione del Sistema di Ipotesi:} Definire esplicitamente $H_0$ e l'alternativa $H_1$ (bilaterale o unilaterale).
    \item \textbf{Scelta della Statistica Test e Livello $\alpha$:} Individuare la statistica pivotale (es. $Z_{\text{test}}$ o $T_{\text{test}}$) e calcolarne la distribuzione esatta sotto l'assunzione che $H_0$ sia vera.
    \item \textbf{Individuazione della Soglia Critica e Calcolo del Valore Osservato:}
    \begin{equation}
        Z_{\text{oss}} = \frac{\overline{x} - \mu_0}{\sigma / \sqrt{n}} \qquad \text{oppure} \qquad T_{\text{oss}} = \frac{\overline{x} - \mu_0}{s / \sqrt{n}}
    \end{equation}
    \item \textbf{Regola di Decisione e Conclusione Statistica:}
    \begin{equation}
        \begin{cases}
            T_{\text{oss}} \in \mathcal{R} \implies \textbf{Rifiuto } H_0 \text{ al livello } \alpha \text{ a favore di } H_1 \text{ (evidenza significativa);} \\
            T_{\text{oss}} \notin \mathcal{R} \implies \textbf{Non Rifiuto } H_0 \text{ (dati compatibili con } H_0 \text{ a livello } \alpha).
        \end{cases}
    \end{equation}
\end{enumerate}
\end{metodo}

% ==============================================================================
% SEZIONE 11.4: IL CONCETTO DI P-VALUE
% ==============================================================================
\section{Il Concetto di $p$-value (Livello di Significatività Osservato)}
\label{sec:p_value_teoria}

Il confronto rigido tra statistica test e soglia critica prefissata $\alpha$ comunica una decisione binaria (rifiuto / non rifiuto), ma non esprime quanto la decisione sia schiacciante o marginale. Per questa ragione l'analisi statistica contemporanea privilegia il **$p$-value**.

\begin{definizione}{$p$-value (Livello di Significatività Osservato)}{def_p_value}
Il \textbf{$p$-value} è la probabilità, calcolata assumendo rigorosamente vera l'ipotesi nulla $H_0$, di osservare un valore della statistica test uguale o ancor più sfavorevole ad $H_0$ di quello effettivamente calcolato sui dati campionari:
\begin{equation}
    p\text{-value} = \begin{cases}
        2 \cdot \left(1 - \Phi(|z_{\text{oss}}|)\right), & \text{per test bilaterale } (H_1: \mu \neq \mu_0) \\
        1 - \Phi(z_{\text{oss}}), & \text{per test unilaterale destro } (H_1: \mu > \mu_0) \\
        \Phi(z_{\text{oss}}), & \text{per test unilaterale sinistro } (H_1: \mu < \mu_0)
    \end{cases}
\end{equation}
\end{definizione}

\begin{figure}[htbp]
    \centering
    \begin{tikzpicture}[scale=0.85, >=Stealth]
    \draw[->, thick, navyblue] (-4.5,0) -- (4.5,0) node[right, font=\footnotesize] {$z$};
    \draw[->, thick, navyblue] (0,-0.2) -- (0,3.8) node[above, font=\footnotesize] {$f(z \mid H_0)$};

    % Curva H0
    \draw[domain=-4.0:4.0, smooth, variable=\x, very thick, navyblue]
        plot ({\x}, {3.5*exp(-\x*\x/2)});

    % Valore osservato z_oss = 2.15
    \fill[warmamber!60, domain=2.15:4.0, variable=\x]
        (2.15,0) -- plot ({\x}, {3.5*exp(-\x*\x/2)}) -- (4.0,0) -- cycle;
    \fill[warmamber!60, domain=-4.0:-2.15, variable=\x]
        (-4.0,0) -- plot ({\x}, {3.5*exp(-\x*\x/2)}) -- (-2.15,0) -- cycle;

    \draw[ultra thick, dashed, warmamber!90!black] (2.15,0) -- (2.15, {3.5*exp(-2.15*2.15/2)});
    \draw[ultra thick, dashed, warmamber!90!black] (-2.15,0) -- (-2.15, {3.5*exp(-2.15*2.15/2)});

    \node[below, font=\scriptsize\bfseries, text=warmamber!90!black] at (2.15,0) {$+z_{\text{oss}}$};
    \node[below, font=\scriptsize\bfseries, text=warmamber!90!black] at (-2.15,0) {$-z_{\text{oss}}$};

    \node[font=\scriptsize\bfseries, text=warmamber!90!black, align=center] at (3.2,1.2) {Area $= \dfrac{p\text{-value}}{2}$};
    \node[font=\scriptsize\bfseries, text=warmamber!90!black, align=center] at (-3.2,1.2) {Area $= \dfrac{p\text{-value}}{2}$};

    \node[font=\footnotesize\itshape, text=navyblue, align=center] at (0,-1.0) {
        $p\text{-value} = 2 \cdot (1 - \Phi(|z_{\text{oss}}|))$. Se $p\text{-value} \le \alpha \implies$ Rifiuto $H_0$.
    };
\end{tikzpicture}

    \caption{Significato geometrico del $p$-value in un test bilaterale. Il $p$-value corrisponde all'area cumulata delle code estreme esterne a $\pm |z_{\text{oss}}|$. Se l'area è inferiore ad $\alpha$, il risultato è talmente anomalo sotto $H_0$ da imporne il rifiuto.}
    \label{fig:p_value_interpretazione}
\end{figure}

\begin{teorema}{Regola di Decisione Universale basata sul $p$-value}{thm_regola_pvalue}
La regola decisionale universale stabilisce:
\begin{equation}
    \begin{cases}
        p\text{-value} \le \alpha \implies \textbf{Rifiuto } H_0 \text{ (il dato è statisticamente significativo a livello } \alpha); \\
        p\text{-value} > \alpha \implies \textbf{Non Rifiuto } H_0 \text{ (evidenza non sufficiente a respingere } H_0).
    \end{cases}
\end{equation}
\end{teorema}

% ==============================================================================
% SEZIONE 11.5: DUALITÀ TRA TEST E INTERVALLI DI CONFIDENZA
% ==============================================================================
\section{Teorema di Dualità tra Test Bilaterali e Intervalli di Confidenza}
\label{sec:dualita_test_ic}

Esiste un profondo legame strutturale e algebrico tra l'inferenza per intervalli e la verifica delle ipotesi.

\begin{teorema}{Dualità Test-Intervalli di Confidenza}{thm_dualita}
Sia $\operatorname{IC}_{1-\alpha}(\theta) = [L(\mathbf{X}), U(\mathbf{X})]$ l'intervallo di confidenza bilaterale al livello $(1-\alpha) \cdot 100\%$ per il parametro $\theta$.
Allora il test d'ipotesi bilaterale $H_0: \theta = \theta_0$ contro $H_1: \theta \neq \theta_0$ al livello di significatività $\alpha$ produce la medesima decisione:
\begin{equation}
    \textbf{Non Rifiuto di } H_0 \text{ al livello } \alpha \iff \theta_0 \in \operatorname{IC}_{1-\alpha}(\theta)
\end{equation}
Equivalentemente, si rifiuta $H_0$ se e solo se il valore ipotizzato $\theta_0$ cade al di fuori dell'intervallo di confidenza: $\theta_0 \notin [L(\mathbf{x}), U(\mathbf{x})]$.
\end{teorema}
